\section{Certificate Spine}
\label{sec:certificates}

The paper-facing certificate table is maintained in \texttt{paper/SOURCE\_MAP.md}.
Before public export, every displayed number in the manuscript must be tied to
one source artifact and one regression test.  The local test count is deliberately
not hardcoded in the mathematical text; it is a build artifact recorded by the
current regression run.

The key public-safe statement is that the B3 matrix is finite-certified, while
the symbolic derivation of the count \(60\) remains open.  The count-collapse
certificate reduces that debt to the finite endpoint quotient kernel
\[
  \delta=0,4,6 \mapsto 36,16,7
\]
and to the B3 row-degree pattern
\[
  1,1,1,2,2,1,4,3.
\]
The intrinsic column-classification certificate closes the column side by proving that all eight B3 profile columns are
intrinsic to \((S,B)\).  The defect-6 realizability-obstruction certificate does not close the realization side; it proves
that the defect-6 residue is an irreducible integer-magic cut, not a consequence
of \(\F_2\)-completability or linear symmetry.

\begin{certprop}[Reproducibility artifact table]
The manuscript depends on the following certified artifacts.  Exact file names
are kept in \texttt{paper/SOURCE\_MAP.md}; this table uses only reproducibility
handles and mathematical block names.
\begin{center}
\scriptsize
\begin{tabular}{lll}
\toprule
Claim block & Reproducibility handle & Status\\
\midrule
endpoint split & endpoint split certificate & certified\\
defect dichotomy & first-principles defect dichotomy & structural\\
error parity & error-fiber parity bridge & structural\\
B3 matrix & support-action matrix summary & finite\\
B3 intrinsic columns & intrinsic column classification & finite\\
B3 defect-6 residue & defect-6 realizability obstruction & finite\\
endpoint transport & endpoint transport package & finite\\
\(\TK\) freeness & global free-transport theorem & finite\\
B2 package & B2 endpoint boundary package & finite\\
count collapse & count-collapse kernel geometry & finite\\
\bottomrule
\end{tabular}
\end{center}
\end{certprop}

\begin{remark}[Attribution and proof source]
The row/column nim-sum property is positioned against the classical order-four
literature, including Ollerenshaw--Bondi and Ward.  Until a precise external
bibliographic statement is pinned down, the manuscript treats the local Step4
derivation as the proof source.
\end{remark}